\chapter{第六集 \quad 殷殇}

\section{正文}

传说中，女子简狄水中行浴，见玄鸟堕其卵，简狄取而吞之，因孕生契，契长大后辅佐大禹治水有功，受封于商，这以后便有了商族。商代甲骨文中，商王祖先亥的名号中冠以鸟形，这正是商人对玄鸟感生神话的崇信。从太行山中，到冀南、豫北平原，历代先祖经历几多艰险、几多征战，才由弱变强。作为玄鸟子孙，成汤得天命，自奋蹄，约公元前1600年，在山西南部的鸣条之战中，夏王帝桀被成汤率领的商军彻底击溃。而洛水之阳的斟鄩之地，辉煌一时的二里头都城已现衰败之象。大型礼仪和政治性工程因遭破坏而废弃，来自海岱地区的岳石文化，以及被认为与商人先祖有关的下七垣文化人群，大举进入城中，商人还在二里头都邑东北约6公里外建起了偃师商城。

城门狭窄，易守难攻，加之城内有大量储粮囷仓和储物府库，厉兵秣马的气息扑面而来，也许商军就屯驻在这座固若金汤的城池里，监视着不远处那二里头城中的夏遗民。与此同时，古济水之南，一座占地约25平方公里的巨大都城在今天的河南郑州拔地而起。它由内城和外郭组成，其间错落分布着宫城、居住区、作坊区和墓葬区。城垣经历了3000余年的风吹雨蚀、兵火燹灾，至今仍保留着高出地面4~5米、宽20~30米的规模。商人，将这座立国之都称为“亳”，成汤正是在这里统治初生的商王朝。商代早期，逢连年大旱，百姓陷于困厄之中，成汤决定亲赴桑林祈雨，他并未遵从杀人献祭的惯例，而是剪下自己的头发与指甲，放于柴堆之上，作为身体的一部分代己为祭。

兵力强盛，人心归附，商很快成为大邑。《诗经·殷武》中追忆道，“昔有成汤，自彼氐羌，莫敢不来享，莫敢不来王，曰商是常。”以郑州商城和偃师商城为核心，商王朝开始构建覆盖黄河、长江流域的庞大城池网络，一边沿用或改造夏人留下的城池，一边按照新的理念建设新城。中条山下，垣曲商城衔山抱水。长江北岸，盘龙城高大崔嵬。盘龙城的布局，与郑州偃师商城类似，主要由宫殿和手工业作坊区等构成，埋葬风俗和出土的青铜礼器也几乎与中原地区保持一致。盘龙城中的宫殿基址，虽然仅有郑州商城的十分之一，但在当时的长江地区，规格却无有出其右者。这座城，是商人在南方疆域的政治和经济中心，或许，也是控制南方珍贵资源的最佳中转地。

几乎与郑州商城建起同时，冶炼青铜的烈火，就已经照亮了郑州南关外的天空。此处的大型铸铜作坊，生产规模惊人，工艺精湛，不仅制造小件的青铜兵器和工具，还能铸造鼎、鬲、爵、觚等青铜礼器。商人先祖曾生活在北方，很早便接触到经草原传来的青铜冶铸术，后来，更是吸收了二里头的影响，最终发展出那个时代最尖端的大型青铜器铸造技术。此时，熔炉里的铜锡逐渐熔化为合金液体，司炉们敲碎浇道口的炉壁，金色的溶液奔流而出，注入浇铸孔，待其冷却后打碎外范，则大鼎可成。青铜，是商王朝至为重要的政治资源，也象征着这个正冉冉升起的新王朝那举世无双的实力。为了治理越来越广大的疆土，商王朝采取了内外服分而治之的统治形式。

所谓“内服”，就是商王直辖之地，委派官员治理；而“外服”，则是王室贵族的领地，或吸纳臣服的方国，拱卫在王畿四周，既是防御和缓冲地带，也是对外扩张的前线。早商王朝的征途似乎一路奏凯，然而，历史的伏笔此刻再现，改变了时局。建城约200年之后，郑州商城逐渐走向衰落，宫室少有人迹，作坊不事生产。考古工作者在城门附近，发现了这一时期的三处大型青铜器窖藏，在深达4~6米的窖坑中，掺杂朱砂或铺垫木板，藏有不同数量和种类的青铜礼器。在那个时代，少有一座城能够长久地作为都邑，王朝需要寻找新的、能维续庞大人口生存的地方。或许在离开旧都前，商人举行了最后的庄严祭祀，将青铜重器瘗埋于地下，为的是祈求祖先和神灵的护佑。

当郑州商城衰落之际，距其20公里外的索须河畔，却生长起一个大型的中心都邑——小双桥，城中那大型的礼仪性建筑、大规模的祭祀活动、工艺精湛的青铜重器，都彰显着这并非一座普通的城址，而可能是商王朝的第二座都城——隞。在这座城内，除了商人的活动痕迹之外，竟还出现了一些具有山东海岱地区文化特征的器物，城中宏伟的高台祭坛，曾举行过比郑州商城中规模更大的祭祀活动；数个祭祀坑内，都填埋着大量非正常死亡的人骨。此时的商人，似已忘却了当年成汤以己身代人牲的苦心。仲丁统治期间，迁都至隞，此时，商人与东夷的关系逐渐紧张，仲丁征战告捷后，带回大量东夷俘虏，以他们为人牲举行了献祭仪式。不久后，一支东夷军队却突袭隞都，杀死了大量来不及逃走的商人。

几番鏖战后，商王朝虽然夺回都城，却也元气大伤。更不利的是，商人王位的传承次序向不固定，极易造成不稳定的局面。仲丁以后，王室贵族们为了王位同室操戈，由此引发的混乱波及九代商王。由于各种因素的叠加，商人不得不放弃旧都，离开中原，踏上了在黄河两岸间来回迁徙的漫长路途。直到距今约3300年前，商王盘庚迁都至殷墟，局势才自此安定。大邑商，这气势恢宏的晚商王都，继续着“四方之极”的构建。西北方的高冈上，商王陵与2800多座祭祀坑俯临俗世。王都内，宫殿和宗庙被环护在洹河和壕沟共同组成的森严屏翳中。西南角，是祭祀王朝土地之神的“社”，东南角是祭祀商人远祖的“宗”。社与宗之北，是为外朝，典礼和仪式活动都在此处举行。

再北则是内寝，也就是商王的内宫，始建于武丁时期，沿用到帝乙、帝辛时期。寝宫的长案之上，摆着几件青铜盥洗之器，侍女向铜盂里放入香茅，再将青铜鼎中的热水舀入盂中，一时满室馨香。铜盘盛满清水，可以调节盂中的水温。商王武丁面色凝重，进行着占卜前的准备，侍女将布浸在盂中，拧至半干，为武丁精心擦拭身体。室内火烛闪烁，映照着武丁健壮的身躯，那一道道伤痕是无数次征战的痕迹，侍女手握陶磢抚摩皮肤，以去除污垢。武丁的心事，是妻子妇好分娩过程的吉凶。妇好是他的配偶，是孩子的母亲，也是朝堂上的股肱。她主持祭祀，领兵征伐，于国于家，不可或缺。邦畿千里，日理万机，妇好的事总是让武丁牵挂：

她每有牙痛、鼻患、腹痛或是偶感风寒，武丁都会郑重地占卜吉凶，祭祀消灾。数年前，他为将要临盆的妇好卜问吉日，31天后的甲寅日，妇好分娩，却偏偏不是求得的吉日，这一次，他不愿再有任何闪失。卸下戎装的妇好仍是个爱美的女子，她最喜欢玉，玉牛、玉熊、玉蛙、玉羊首、玉蝉、玉虎、玉人，这些精雕细琢的美玉，凝聚着天地的灵气。玉镯、玉串饰和玛瑙串饰，还有精致的骨笄，摆满了她的首饰盒。她常常对镜梳妆，即使对王室来说，铜镜也是稀罕之物，身份尊贵如妇好，方能拥有。她用玉梳精心梳理并挽起长发，再从众多的发饰中，挑出一只夔首骨笄，紧紧固定住发髻。今夜皓月当空，商王与王后，一起来至早已备好的祭场中进行祷告。

这钺，是妇好征战时号令军旅的统帅标志，也是仪式中驱祟的法器，钺未装柄时便重达9公斤。钺身上饰有侧身对立、张口咆哮的双虎，虎口相对间嵌有一颗人首，虎身后各饰一条卷尾之龙，虎爪下是巨大的獠牙。商人崇信神祇，通过这样的仪式，能得神灵与祖先的护佑，顺利产子，延续商祚。此刻的密室中央，炭火熊熊，贞人宾拿起卜骨，在炭火上轻轻烘烤，身旁的矮几上放着整治、钻凿和刻画甲骨的工具。贞人，不但是占卜活动的具体操作者，还承担着巫的职责。少数掌管占卜机构的卜官，是商王之侧的高级贵族。武丁左手扶住卜骨，右手持炭棒灼烧着钻槽，通红的棒头神秘地闪烁，不久后，裂缝在钻槽边绽开，发出脆响。

武丁屏息细看卜骨上裂开的兆纹，脸上渐渐露出喜色，这次占卜所得是大吉之兆。“丁酉卜，宾贞：妇好有受生？”“王占曰：吉！”“其有受生！”

宾拿起青铜刻刀，开始在龟甲上契刻。传说仓颉造出文字之时，天雨粟、鬼夜哭，天地间同为这亘古未有的创造而震撼。文字之于古人，如劈入混沌世界的雷电，惊醒并照亮了蛰伏的文明。自此以后，人类拥有了能够世代累积和叠加的智慧与情怀，万物灵长之文明永难磨灭。从武丁时代开始，中国拥有了迄今可以确认的最早的成熟文字系统。象形文字，是商朝的祖先留给我们最宝贵的遗产。至今，这方正汉字的一笔一画间，仍存续着商人眼中那日升月落、风雨雷电的万象，古今攸同，仍维系着这三千多年来每一颗以汉字启蒙的心灵，无问西东。

武丁时期，开始实行“周祭”制度，也就是轮流祭祀自先祖上甲以来的历代先公、先王和王后，周而复始，一年完成一个祭祀周期。而升入天宫的历代先王们，此刻正陪伴在“帝”的身边，“帝”是至高之神，可赐予丰收，护佑凯旋，唯有在世商王才具备与祖先和“帝”沟通的能力。通过祭祀、占卜，获取“帝”的旨意和世间人事的预兆，祖先的遗骨垒出了信仰的神殿，而商王，正是唯一的大祭司。丰盛的食物来自于天赐与人功，也需以青铜盛装，才显出在人间的尊崇和对神明的敬重。巨大的方鼎，可烹煮大型祭牲，造型浑厚而庄重，如天地间最均衡和稳定的存在。簋，盛满煮熟的谷物；豆，盛放肉类或调料；甗，是复合器物，下为鬲，上为甑，可实现蒸煮一体的功能。

佳肴粢盛，盛以吉金重器，神祇与祖先想必能尽情享用。商人嗜酒，酒也是祭祀时不可或缺的祭品，它有助于让司祭者进入一种仿若通灵的状态。因此在青铜器中，酒器便成为最重要的一种器类。盛酒的瓿、彝、罍、尊、盉、卣、壶、觥；温酒的爵、角、鬲；饮酒的觚、觯，玉醴香醪，配以金铜炫曜，这来自人间的美酒芬芳想必能直达云霄。这日，武丁接到西边的周族被降服的消息，尽管当时生活在关中平原的这支小族实在势微力弱、不值一提，但恰逢大祭前夕，武丁觉得是吉兆，仍十分高兴。

每逢献祭，武丁便饮尽杯中酒，他注视着眼前那尊金色的大鼎，燎柴之烟氤氲缭绕，神光离合间，鼎上的蟠龙升腾，神鸟飞起，那双目炯炯的饕餮幻化为庞然巨物，凌驾于上空，恍惚间又只余缕缕青烟，与天上云气若混作一体，不知何处飞来的一只锦雉，竟立于鼎上，引吭清啼。武丁大感惶恐，生怕祭仪中这样一个偶发事件会成为不祥之兆，身旁的臣子祖己却劝解商王，只要勤修政事，便无可畏惧。

商代晚期的青铜器，各种器形已臻于成熟之境，与稳重简洁的造型形成对比的是纹饰的繁缛。铺满整器的地纹中，又以更粗重的线条与凹凸，勾勒出各种具有丰富意象的纹样。

狞厉的饕餮、蜷曲的夔龙、昂首的凤鸟、重生的蜕蝉，这些复杂纹饰熔铸于圆弧的容器表面，却无一处废笔或错漏，非高超至极的工艺和设计能力不能达成。商人通过精湛的青铜重器，营造了一种光怪陆离、神秘莫测的氛围，这也是商人浩荡而深窈的精神世界。神灵漫天、异兽满地，人之在其中，惟有商王可以沟通天地。而这些集稀缺资源、复杂工艺和神秘知识于一体的青铜礼器，其组织生产、使用方法和器用制度，也都一概掌握在商王的手中。独占神权，只实现了精神层面的统摄，这尚不足以使商的统治者高枕无忧。唯有配合以强大的军事实力，才能构成王位之畔的坚固藩篱。

商的军队组织周密而完备，装备有最精良的杀伐利器，钺是砍杀之器，也是军权和王权的象征；戈宜钩杀，矛宜刺杀，强弓劲矢可远程伏射，甲胄能护身御敌。狩猎，是日常练兵的一种方式，远处一头野牛出现，武丁令战车驱逐，不料，身后的小臣叶所驾之车却撞将上来，两驾车都毁坏，王子子央也堕车受伤。祭祀狩猎涂朱牛骨刻辞，记录了这起追尾事故，在写到“车”字时，利用象形文字的优势，极形象地展示了车祸现场，一“车”车轴断裂，一“车”倒翻。

距今约4000年前，乌拉尔山地区出现了世界上最早的马拉车，中国最早的成熟马车则出现于晚商的武丁时期，形制与欧亚草原的马车相似，匹配刀、弓形器、马策等工具组，表明其驾驭方式也与草原马车近似，暗示这些马车是草原地区传入，而非本土起源。

但以成套铜饰件装饰马车的做法则是商人的发明。马车的出现改变了战争形态，在此后的900年间，车战成为古中国大规模战争的主要形式，千乘之国是强国的象征，万乘之尊是帝王的称号。战争于胜者而言是荣耀，却也不可避免会带来伤亡。这位年轻的商军士兵被敌人箭镞射入头骨，死在了战场之上。即便在战场上侥幸苟活，一旦被俘，却可能面临着更加悲惨的命运，作为人牲献祭，是大多数战俘最终的结局。卜辞记载，武丁时期用于祭祀的人牲数量可达9021人，最大规模的一次足足献祭了500名人牲。殷墟曾出土过两片刻有文字的人头骨，其中一片头骨上刻写的内容是“方伯用”，意思是用“方伯”来献祭；另一篇则刻了“方伯，祖乙伐”，也就是说，用“方伯”来献祭商王祖乙。

方伯，是商人对方国首领的称呼。商人将这两位怀有不臣之心的一方之主制成了刻字的祭品，以此献给那长随天地之侧的祖先。“大邑商”，是商人对自己这巍峨国都的称呼，宫殿、宗庙与王陵，在世俗政治和宗教神权体系里，都有着至高无上的地位，统摄着整个王朝。以殷墟为中心，北起沙丘、南到朝歌、西至太行山、东抵古黄河，这便是商的王畿所在。水渠和道路由这中心生长出去，如虬须盘根般蔓延至商的每个角落，由王畿向外延伸数百公里。侯、甸、男、卫等各级贵族的领地，构成商的藩屏，方国势力交错其间。黄昏时分的豫北平原，北行的牛车队伍，正将四方奉献的贡品载向王都，其中不仅有谷物、象牙、龟甲，还有海贝、铜锭，甚至是鲸鱼骨骼。

另一辆双马轻车在路上奔驰，驭手是长族首领亚长的属下，亚长奉商王之命出征，已领军与东夷人激战月余，眼见胜利在望，亚长令驭手返回王都，向商王汇报战况。此刻天色近晚，不远处便是“羁”，也就是商人的驿站，按照原定的计划，在此休息一夜，明日一早起身赶路，午间便能回到王都。驭手念及将能见到分别许久的家人，心中欢喜，亚长献给商王的战利品，由驿卒代为看管，这包裹里却是他自己攒下来留给家人的礼物，给妻子的玉饰、给儿子的青铜刀、给女儿的海螺。他检视了一番，确定这些宝物并未因路途的颠簸而毁损，才放心地珍重藏好。驿站外忽然传来一阵急促的马蹄声，一位骑士推门进屋。

驭手认得，这是亚长的侍卫，他的突然出现让驭手心中十分不安。果然，侍卫带来了噩耗，在驭手出发后不久，本以为稳操胜券的亚长竟血染疆场。侍卫转告驭手不必再赶回王都，因为护送遗体的人马，将于三日后启程返回大邑商，驭手需在驿站中等候，而侍卫则要星夜兼程赶回王都报讯。驭手取出那个珍藏已久的包裹交给侍卫，嘱托他一定要带给自己的家人。驭手心中悲痛，他知道这样好的夜色，自己不能再看几回了。按照商人的习俗和礼仪，亚长一死，意味着作为亚长家臣的他，势必要殉葬于地下。迎回亚长遗体的家人，正沉浸于一片哀痛中。

亚长的遗体已经清洗干净，敌人集中袭击了他的身体中部，左大腿根部正面被锐斧劈砍，他同时受到来自左后方的利戈的袭击，动脉被割破，盆骨右侧被长矛刺穿，左臂留下至少3处锐器砍伤痕迹，左肋被锐器击伤。亚长伤痕累累的身体，将被花椒完全覆盖，这是亚长家乡特有的防腐方法，外面再裹以麻布和锦缎，面朝下放入棺中，这样的俯身葬也是家乡的习俗。深夜，乡村的茅草屋中，驭手的妻子仍等待着离人归来，这战火再起的时代，或许没有消息就是最好的消息。马车，本是身份地位的象征，乘坐马车者，通常为王公贵族，他们死后，马车将作为陪葬品埋于墓内或其附近，马匹和驭手常常被一并埋入。冬日阴霾的天空下，西风萧瑟，犬鸣马啸。

驭手跪坐在商王为亚长择定的墓旁，和他一起的，还有另外14位将一同殉葬的家臣，其中也有那夜报讯的侍卫，那个面容稚嫩的年轻人。有人牵来了要殉葬的狗；临行前，驭手饮了些酒，这是他人生的最后一杯酒。

殷墟王宫区的东南侧，花园庄东地的54号大墓中，墓主是35岁左右的长族首领亚长，与他一同下葬的还有15位殉人和15条殉狗，各类随葬品约1800余件，其中方鬲、方鼎、牛尊等，是高等级人群方能使用的器物。青铜礼器，尤其是觚和爵的数量，是社会等级的标志。妇好墓出土53觚40爵，亚长墓随葬9觚9爵，数量之别，正是等级之分。商人构建的这套以青铜器为载体的礼制系统，将祭祀礼器制度与社会等级制度相结合，并以特定器类的数量展现出来。

这样的礼制设计，直观而明确，覆盖面和接受度都很广，从而成为商代构建强大国家体制的重要内容。

对亚长牙齿的锶同位素检测结果显示，他的幼年并非在殷墟度过。他所属的长族，封地在今河南东南部的鹿邑，毗邻东夷和淮夷，是商的东南门户。为了争夺土地，控制海盐，商与东夷交战不休。济南的大辛庄遗址，是一处商王朝深楔入东夷腹地的军事据点，出土了与商王朝同样的青铜礼器组合，还发现了刻辞卜甲，这是晚商时期，殷墟以外罕有的发现了甲骨文的地方，无论是字形、文法，还是甲骨修正、钻凿形态，都与安阳殷墟卜辞属于同一系统。山东青州，则是另一处两军交锋的前沿。苏埠屯1号墓的墓主，是“亚醜”族的首领，采用了殉人、腰坑殉狗等典型的商人葬俗。

带“亚醜”徽记的铜钺，表明他拥有商王赋予的军权，为商镇守东土。

那么，遥远的南土，商人又该如何控制？青年是商的王族，受商王之命，巡视长江沿岸形势，筹划治理南土的良策。早先镇守南土的盘龙城早已凋敝，王朝疲于应付与东夷、西北的舌方（吾方是殷墟出土卜辞中见到的一个商代的方国部落）、鬼方、羌人之间的冲突，久已疏于南土的治理。青年望向江面，回想起在赣江中游所见的奇特祭器，那狰狞的猛虎躯体硕大，呈伏蹲欲纵之态，虎周身遍饰花纹，背上伏卧一只小鸟，宛若猛虎的驾驭者。头生双角的人面，神与人似融合一体，面容神秘诡异，充满威严。青年打开一幅帛画，所绘青铜礼器，正属于湘江中游的方国。四羊方尊、人面方鼎，虽与商器相类，但形象生动，别具特色。

商族青年逆流而上，进入三峡，两岸群山掩映，江水变得湍急起来。此刻，古蜀国神殿的密室中烛火通明，工匠正在打着一副黄金面具。方面大耳，直鼻阔口，三角形眼孔，幽暗深邃。经由宽阔的江河，抵达终古的密林，商族青年似被感召，径直踏入这古蜀国的神秘地界。硕大的太阳转动不休，威严的瞋目俯瞰尘世，商族青年恍惚之中，忽然心惊，商王独占了与帝沟通、祭祀先祖的神权。可是若这天下，各地有各地信仰之神呢？王又如何能控制四域？古蜀神殿的大门，正在徐徐敞开，硕大的青铜面具凸目阔口，造型意象上体现出人神合一的特点，或许是古蜀人的祖先神造像，双手举持法器的青铜立人，可能是其神、巫、王三者身份于一体的权威人物。

青铜神树，似从象征“神山”的底座上生出，枝叶蜷曲，奇花绽放，枝桠上立着昂首的神鸟，树侧有逶迤而下的铜龙。这是青铜鸟足神像，最下部为青铜方座罍，罍表可见少量红彩残留，器物中部为青铜神像，纵目、獠牙、额生牛角，以双手为支撑，倒立在青铜罍之上，他头顶一件青铜觚形尊，尊盖上还站着手持龙形器的立人，整器，似在表现鸟族人面的神祇从天而降。这件青铜神坛的镂空台基之上，共有13个小型青铜人像，其中有一人，以跪姿背罍，似为女性；抬架之上，立有一只如犬似马的青铜神兽，在神兽的头与尾之间，跪坐一人，手举山形座，神坛，由四根细体觚支撑，侧立青铜持鸟立人像，圆盘上方，有四个顶尊跪坐人像。

这件器物，构成了一个祭祀活动场景，可能反映了重大仪式活动中的人员分工以及宗教观念。牙璋，是自二里头继承来的玉礼器，传至古蜀国后，往往在牙璋顶端劈出尖嘴，使其酷似鱼形，有时又在鱼嘴中雕出一只振翅之鸟，仿佛鱼鸟化生。这件跪坐小人像，总高不足5厘米，下身着裙，上身赤裸，双手持握一璋，是古蜀人以璋行祭的实物例证。三星堆人，对商礼制系统中，尊、罍等器物的推崇，以及对中原独有的青铜范铸技术的采用，表现出自身与中原王朝的密切联系。四川广汉三星堆古城，可调控环境参数的考古舱内，精细的发掘正在进行。三千多年前，古蜀人可能在仪式中，亲手砸碎和埋葬了这些青铜器。

三千多年后，经考古工作者的双手，它们重见天日，并将在此后被一一修复、渐渐弥合，从而再现出神秘古蜀的宏大时空。三星堆的历史，是见证中国多元一体交融汇聚的缩影，距今4500年开始，西北地区的先民带着彩陶文化，经过川西山区进入成都平原的西北地区，同时抵达的还有经三峡地区而来的长江中游先民。此后漫长的岁月里，在中原和长江流域的青铜文明的共同作用下，才诞生了古蜀国这样一个区域文明。步入王国社会鼎盛时期的中原王朝，此刻已感知到潜在的危机，由他们树立的青铜礼制、文字体系、神权信仰和军事规范，有力维系和影响着辽阔疆域中的各种力量。但与此同时，也培养了自己的挑战者。

长江沿线，东有江西新干大洋洲文化，中有湖南宁乡青铜文化，西有四川广汉三星堆文化，商的南土，已经被这些新兴方国所控制。在西北地区更是强敌四起，陕北一带，卜辞中的“鬼方”，晋西山地，卜辞中的“舌方”，双双屹立在沟壑纵横的黄土高原。被后世称为纣王的帝辛，征服了强大的东夷后，继而又征服了蓝夷，将淮河和长江下游纳入实际控制区内，《史记·殷本纪》中的帝辛，见闻广博，勇力过人，但穷奢极欲，残暴不仁，其中有多少是史实，多少是周人的贬低，不得而知。殷墟考古中，各个时期的墓葬数量以帝辛时最多，各族邑遗址中发现的大型四合院建筑群、带墓道的大型墓葬及车马坑，也多建于帝辛时期，这表明王朝人口的蓬勃增长和国度的空前繁荣。

但帝辛终未能改变商的命运，靡不有初，鲜克有终，成汤时对子民的顾念，并未被后来的商王们继承。商人好问，他们问天、问神、问祖先，凡事他们都要问卜，直到有一天，神不再应答。不过，来自商的祖先，无意中还留下了他们的样子，这尊陶土塑成的、比心还小的脸庞。他是谁？他是为主人殉葬的驭手，他是向命运索取答案的贞人，他是顺着江河漂泊的青铜匠人……他望着我们，仿佛有很多话要说，那是甲骨文也无法言传的、一个人平凡的悲欢离合……而在西北的黄土高原上，出自戎狄之间的周人已经蓄势待发，他们，将带着我们的中国，走向闪烁人性光芒的家国天下。
